\section{Tame symbols}
\label{section-tame-symbol}

\noindent
Consider a Noetherian local ring $(A, \mathfrak m)$ of dimension $1$.
We denote $Q(A)$ the total ring of fractions of $A$, see
Algebra, Example \ref{algebra-example-localize-at-prime}.
The {\it tame symbol} will be a map
$$
\partial_A(-, -) : Q(A)^* \times Q(A)^* \longrightarrow \kappa(\mathfrak m)^*
$$
satisfying the following properties:
\begin{enumerate}
\item $\partial_A(f, gh) = \partial_A(f, g) \partial_A(f, h)$
\label{item-bilinear}
for $f, g, h \in Q(A)^*$,
\item $\partial_A(f, g) \partial_A(g, f) = 1$
\label{item-skew}
for $f, g \in Q(A)^*$,
\item $\partial_A(f, 1 - f) = 1$
\label{item-1-x}
for $f \in Q(A)^*$ such that $1 - f \in Q(A)^*$,
\item $\partial_A(aa', b) = \partial_A(a, b)\partial_A(a', b)$
\label{item-bilinear-better}
and $\partial_A(a, bb') = \partial_A(a, b)\partial_A(a, b')$
for $a, a', b, b' \in A$ nonzerodivisors,
\item $\partial_A(b, b) = (-1)^m$
\label{item-skew-better}
with $m = \text{length}_A(A/bA)$
for $b \in A$ a nonzerodivisor,
\item $\partial_A(u, b) = u^m \bmod \mathfrak m$
\label{item-normalization}
with $m = \text{length}_A(A/bA)$ for $u \in A$ a unit and
$b \in A$ a nonzerodivisor, and
\item
\label{item-1-x-better}
$\partial_A(a, b - a)\partial_A(b, b) = \partial_A(b, b - a)\partial_A(a, b)$
for $a, b \in A$ such that $a, b, b - a$ are nonzerodivisors.
\end{enumerate}
Since it is easier to work with elements of $A$ we will
often think of $\partial_A$ as a map defined on pairs of
nonzerodivisors of $A$ satisfying (\ref{item-bilinear-better}),
(\ref{item-skew-better}), (\ref{item-normalization}),
(\ref{item-1-x-better}). It is an exercise to see that
setting
$$
\partial_A(\frac{a}{b}, \frac{c}{d}) =
\partial_A(a, c) \partial_A(a, d)^{-1} \partial_A(b, c)^{-1} \partial_A(b, d)
$$
we get a well defined map $Q(A)^* \times Q(A)^* \to \kappa(\mathfrak m)^*$
satisfying (\ref{item-bilinear}), (\ref{item-skew}), (\ref{item-1-x})
as well as the other properties.

\medskip\noindent
We do not claim there is a unique map with these properties.
Instead, we will give a recipe for constructing such a map.
Namely, given $a_1, a_2 \in A$ nonzerodivisors, we choose
a ring extension $A \subset B$ and local factorizations
as in Lemma \ref{lemma-prepare-tame-symbol}.
Then we define
\begin{equation}
\label{equation-tame-symbol}
\partial_A(a_1, a_2) = \prod\nolimits_j
\text{Norm}_{\kappa(\mathfrak m_j)/\kappa(\mathfrak m)}
((-1)^{e_{1, j}e_{2, j}}u_{1, j}^{e_{2, j}}u_{2, j}^{-e_{1, j}}
\bmod \mathfrak m_j)^{m_j}
\end{equation}
where $m_j = \text{length}_{B_j}(B_j/\pi_j B_j)$ and the product
is taken over the maximal ideals $\mathfrak m_1, \ldots, \mathfrak m_r$ of $B$.

\begin{lemma}
\label{lemma-well-defined-tame-symbol}
The formula (\ref{equation-tame-symbol}) determines a
well defined element of $\kappa(\mathfrak m)^*$. In other words, the
right hand side does not depend on the choice of the
local factorizations or the choice of $B$.
\end{lemma}

\begin{proof}
Independence of choice of factorizations. Suppose we have
a Noetherian $1$-dimensional local ring $B$, elements $a_1, a_2 \in B$,
and nonzerodivisors $\pi, \theta$ such that we can write
$$
a_1 = u_1 \pi^{e_1} = v_1 \theta^{f_1},\quad
a_2 = u_2 \pi^{e_2} = v_2 \theta^{f_2}
$$
with $e_i, f_i \geq 0$ integers and $u_i, v_i$ units in $B$.
Observe that this implies
$$
a_1^{e_2} = u_1^{e_2}u_2^{-e_1}a_2^{e_1},\quad
a_1^{f_2} = v_1^{f_2}v_2^{-f_1}a_2^{f_1}
$$
On the other hand, setting
$m = \text{length}_B(B/\pi B)$ and $k = \text{length}_B(B/\theta B)$
we find $e_2 m = \text{length}_B(B/a_2 B) = f_2 k$.
Expanding $a_1^{e_2m} = a_1^{f_2 k}$ using the above we find
$$
(u_1^{e_2}u_2^{-e_1})^m =  (v_1^{f_2}v_2^{-f_1})^k
$$
This proves the desired equality up to signs. To see the signs
work out we have to show $me_1e_2$ is even if and only if
$kf_1f_2$ is even. This follows as both $me_2 = kf_2$ and
$me_1 = kf_1$ (same argument as above).

\medskip\noindent
Independence of choice of $B$. Suppose given two extensions
$A \subset B$ and $A \subset B'$ as in Lemma \ref{lemma-prepare-tame-symbol}.
Then
$$
C = (B \otimes_A B')/(\mathfrak m\text{-power torsion})
$$
will be a third one. Thus we may assume we have
$A \subset B \subset C$ and factorizations over the
local rings of $B$ and we have to show that using
the same factorizations over the local rings of $C$
gives the same element of $\kappa(\mathfrak m)$.
By transitivity of norms
(Fields, Lemma \ref{fields-lemma-trace-and-norm-tower})
this comes down to the following problem:
if $B$ is a Noetherian local ring of dimension $1$
and $\pi \in B$ is a nonzerodivisor, then
$$
\lambda^m = \prod \text{Norm}_{\kappa_k/\kappa}(\lambda)^{m_k}
$$
Here we have used the following notation:
(1) $\kappa$ is the residue field of $B$,
(2) $\lambda$ is an element of $\kappa$,
(3) $\mathfrak m_k \subset C$ are the maximal ideals of $C$,
(4) $\kappa_k = \kappa(\mathfrak m_k)$ is the residue field of
$C_k = C_{\mathfrak m_k}$,
(5) $m = \text{length}_B(B/\pi B)$, and
(6) $m_k = \text{length}_{C_k}(C_k/\pi C_k)$.
The displayed equality holds because
$\text{Norm}_{\kappa_k/\kappa}(\lambda) = \lambda^{[\kappa_k : \kappa]}$
as $\lambda \in \kappa$ and because $m = \sum m_k[\kappa_k:\kappa]$.
First, we have $m = \text{length}_B(B/xB) = \text{length}_B(C/\pi C)$
by Lemma \ref{lemma-compare-periodic-lengths} 
and (\ref{equation-multiplicity-coker-ker}).
Finally, we have $\text{length}_B(C/\pi C) = \sum m_k[\kappa_k:\kappa]$
by Algebra, Lemma \ref{algebra-lemma-pushdown-module}.
\end{proof}

\begin{lemma}
\label{lemma-tame-symbol}
The tame symbol (\ref{equation-tame-symbol}) satisfies
(\ref{item-bilinear-better}), (\ref{item-skew-better}),
(\ref{item-normalization}), (\ref{item-1-x-better}) and hence
gives a map $\partial_A : Q(A)^* \times Q(A)^* \to \kappa(\mathfrak m)^*$
satisfying (\ref{item-bilinear}), (\ref{item-skew}), (\ref{item-1-x}).
\end{lemma}

\begin{proof}
Let us prove (\ref{item-bilinear-better}).
Let $a_1, a_2, a_3 \in A$ be nonzerodivisors.
Choose $A \subset B$ as in Lemma \ref{lemma-prepare-tame-symbol}
for $a_1, a_2, a_3$. Then the equality
$$
\partial_A(a_1a_2, a_3) = \partial_A(a_1, a_3) \partial_A(a_2, a_3)
$$
follows from the equality
$$
(-1)^{(e_{1, j} + e_{2, j})e_{3, j}}
(u_{1, j}u_{2, j})^{e_{3, j}}u_{3, j}^{-e_{1, j} - e_{2, j}} =
(-1)^{e_{1, j}e_{3, j}}
u_{1, j}^{e_{3, j}}u_{3, j}^{-e_{1, j}}
(-1)^{e_{2, j}e_{3, j}}
u_{2, j}^{e_{3, j}}u_{3, j}^{-e_{2, j}}
$$
in $B_j$. Properties (\ref{item-skew-better}) and
(\ref{item-normalization}) are equally immediate.

\medskip\noindent
Let us prove (\ref{item-1-x-better}). Let $a_1, a_2, a_1 - a_2 \in A$
be nonzerodivisors and set $a_3 = a_1 - a_2$.
Choose $A \subset B$ as in Lemma \ref{lemma-prepare-tame-symbol}
for $a_1, a_2, a_3$. Then it suffices to show
$$
(-1)^{e_{1, j}e_{2, j} + e_{1, j}e_{3, j} + e_{2, j}e_{3, j} + e_{2, j}}
u_{1, j}^{e_{2, j} - e_{3, j}}
u_{2, j}^{e_{3, j} - e_{1, j}}
u_{3, j}^{e_{1, j} - e_{2, j}} \bmod \mathfrak m_j = 1
$$
This is clear if $e_{1, j} = e_{2, j} = e_{3, j}$.
Say $e_{1, j} > e_{2, j}$. Then we see that $e_{3, j} = e_{2, j}$
because $a_3 = a_1 - a_2$ and we see that $u_{3, j}$
has the same residue class as $-u_{2, j}$. Hence
the formula is true -- the signs work out as well
and this verification is the reason for the choice of signs
in (\ref{equation-tame-symbol}).
The other cases are handled in exactly the same manner.
\end{proof}

\begin{lemma}
\label{lemma-norm-down-tame-symbol}
Let $(A, \mathfrak m)$ be a Noetherian local ring of dimension $1$.
Let $A \subset B$ be a finite ring extension with $B/A$
annihilated by a power of $\mathfrak m$ and $\mathfrak m$ not
an associated prime of $B$.
For $a, b \in A$ nonzerodivisors we have
$$
\partial_A(a, b) = \prod
\text{Norm}_{\kappa(\mathfrak m_j)/\kappa(\mathfrak m)}(\partial_{B_j}(a, b))
$$
where the product is over the maximal ideals $\mathfrak m_j$ of $B$
and $B_j = B_{\mathfrak m_j}$.
\end{lemma}

\begin{proof}
Choose $B_j \subset C_j$ as in
Lemma \ref{lemma-prepare-tame-symbol} for $a, b$.
By Lemma \ref{lemma-glue-at-max} we can choose a finite ring
extension $B \subset C$ with $C_j \cong C_{\mathfrak m_j}$ for all $j$.
Let $\mathfrak m_{j, k} \subset C$ be the maximal ideals of $C$
lying over $\mathfrak m_j$. Let
$$
a = u_{j, k}\pi_{j, k}^{f_{j, k}},\quad
b = v_{j, k}\pi_{j, k}^{g_{j, k}}
$$
be the local factorizations which exist by our choice of
$C_j \cong C_{\mathfrak m_j}$. By definition we have
$$
\partial_A(a, b) = 
\prod\nolimits_{j, k}
\text{Norm}_{\kappa(\mathfrak m_{j, k})/\kappa(\mathfrak m)}
((-1)^{f_{j, k}g_{j, k}}u_{j, k}^{g_{j, k}}v_{j, k}^{-f_{j, k}}
\bmod \mathfrak m_{j, k})^{m_{j, k}}
$$
and
$$
\partial_{B_j}(a, b) = 
\prod\nolimits_k
\text{Norm}_{\kappa(\mathfrak m_{j, k})/\kappa(\mathfrak m_j)}
((-1)^{f_{j, k}g_{j, k}}u_{j, k}^{g_{j, k}}v_{j, k}^{-f_{j, k}}
\bmod \mathfrak m_{j, k})^{m_{j, k}}
$$
The result follows by transitivity of norms
for $\kappa(\mathfrak m_{j, k})/\kappa(\mathfrak m_j)/\kappa(\mathfrak m)$, see
Fields, Lemma \ref{fields-lemma-trace-and-norm-tower}.
\end{proof}

\begin{lemma}
\label{lemma-tame-symbol-formally-smooth}
Let $(A, \mathfrak m, \kappa) \to (A', \mathfrak m', \kappa')$
be a local homomorphism of Noetherian local rings. Assume $A \to A'$
is flat and $\dim(A) = \dim(A') = 1$. Set
$m = \text{length}_{A'}(A'/\mathfrak mA')$.
For $a_1, a_2 \in A$ nonzerodivisors
$\partial_A(a_1, a_2)^m$ maps to $\partial_{A'}(a_1, a_2)$
via $\kappa \to \kappa'$.
\end{lemma}

\begin{proof}
If $a_1, a_2$ are both units, then $\partial_A(a_1, a_2) = 1$
and $\partial_{A'}(a_1, a_2) = 1$ and the result is true.
If not, then we can choose a ring extension $A \subset B$ and
local factorizations as in Lemma \ref{lemma-prepare-tame-symbol}.
Denote $\mathfrak m_1, \ldots, \mathfrak m_m$
be the maximal ideals of $B$. Let $\mathfrak m_1, \ldots, \mathfrak m_m$
be the maximal ideals of $B$ with residue fields $\kappa_1, \ldots, \kappa_m$.
For each $j \in \{1, \ldots, m\}$ denote $\pi_j \in B_j = B_{\mathfrak m_j}$
a nonzerodivisor such that we have factorizations
$a_i = u_{i, j}\pi_j^{e_{i, j}}$ as in the lemma.
By definition we have
$$
\partial_A(a_1, a_2) = \prod\nolimits_j
\text{Norm}_{\kappa_j/\kappa}
((-1)^{e_{1, j}e_{2, j}}u_{1, j}^{e_{2, j}}u_{2, j}^{-e_{1, j}}
\bmod \mathfrak m_j)^{m_j}
$$
where $m_j = \text{length}_{B_j}(B_j/\pi_j B_j)$.

\medskip\noindent
Set $B' = A' \otimes_A B$. Since $A'$ is flat over $A$ we see
that $A' \subset B'$ is a ring extension with $B'/A'$ annihilated
by a power of $\mathfrak m'$. Let
$$
\mathfrak m'_{j, l},\quad l = 1, \ldots, n_j
$$
be the maximal ideals of $B'$ lying over $\mathfrak m_j$. Denote
$\kappa'_{j, l}$ the residue field of $\mathfrak m'_{j, l}$. Denote
$B'_{j, l}$ the localization of $B'$ at $\mathfrak m'_{j, l}$.
As factorizations of $a_1$ and $a_2$ in $B'_{j, l}$
we use the image of the factorizations
$a_i = u_{i, j} \pi_j^{e_{i, j}}$ given to us in $B_j$.
By definition we have
$$
\partial_{A'}(a_1, a_2) = \prod\nolimits_{j, l}
\text{Norm}_{\kappa'_{j, l}/\kappa'}
((-1)^{e_{1, j}e_{2, j}}u_{1, j}^{e_{2, j}}u_{2, j}^{-e_{1, j}}
\bmod \mathfrak m'_{j, l})^{m'_{j, l}}
$$
where $m'_{j, l} = \text{length}_{B'_{j, l}}(B'_{j, l}/\pi_j B'_{j, l})$.

\medskip\noindent
Comparing the formulae above we see that it suffices to show that
for each $j$ and for any unit $u \in B_j$ we have
\begin{equation}
\label{equation-to-prove}
\left(\text{Norm}_{\kappa_j/\kappa}(u \bmod \mathfrak m_j)^{m_j}\right)^m
=
\prod\nolimits_l
\text{Norm}_{\kappa'_{j, l}/\kappa'}(u \bmod \mathfrak m'_{j, l})^{m'_{j, l}}
\end{equation}
in $\kappa'$. We are going to use the construction of determinants of
endomorphisms of finite length modules in
More on Algebra, Section \ref{more-algebra-section-determinants-finite-length}
to prove this. Set $M = B_j/\pi_j B_j$. By
More on Algebra, Lemma \ref{more-algebra-lemma-multiplication} we have
$$
\text{Norm}_{\kappa_j/\kappa}(u \bmod \mathfrak m_j)^{m_j} =
\det\nolimits_\kappa(u : M \to M)
$$
Thus, by
More on Algebra, Lemma \ref{more-algebra-lemma-flat-base-change-det},
the left hand side of (\ref{equation-to-prove}) is equal to
$\det_{\kappa'}(u : M \otimes_A A' \to M \otimes_A A')$.
We have an isomorphism
$$
M \otimes_A A' = (B_j/\pi_j B_j) \otimes_A A' =
\bigoplus\nolimits_l B'_{j, l}/\pi_j B'_{j, l}
$$
of $A'$-modules. Setting $M'_l = B'_{j, l}/\pi_j B'_{j, l}$ we see that
$\text{Norm}_{\kappa'_{j, l}/\kappa'}(u \bmod \mathfrak m'_{j, l})^{m'_{j, l}}
= \det_{\kappa'}(u_j : M'_l \to M'_l)$ by
More on Algebra, Lemma \ref{more-algebra-lemma-multiplication} again.
Hence (\ref{equation-to-prove}) holds by multiplicativity of the determinant
construction, see More on Algebra, Lemma \ref{more-algebra-lemma-ses}.
\end{proof}






